% ---------------------------------------------------------------------------
% Lista de abreviaturas (norma 3.11)
% ---------------------------------------------------------------------------
% Una fila por cada sigla usada en el informe, en orden alfabético.

\seccionSinNumerar{Lista de abreviaturas}

\begin{longtable}{@{}p{0.18\textwidth}p{0.78\textwidth}@{}}
\toprule
\textbf{Sigla} & \textbf{Significado} \\
\midrule
\endfirsthead
\toprule
\textbf{Sigla} & \textbf{Significado} \\
\midrule
\endhead
\bottomrule
\endfoot

ACS      & \textit{Auto Configuration Server} --- servidor de autoconfiguración remota de
           equipos de cliente bajo TR-069/TR-369 \\
API      & \textit{Application Programming Interface} --- interfaz de programación de
           aplicaciones \\
ATT\&CK  & \textit{Adversarial Tactics, Techniques, and Common Knowledge} --- matriz de
           tácticas y técnicas adversarias de MITRE \\
CLI      & \textit{Command-Line Interface} --- interfaz de línea de comandos \\
CPU      & \textit{Central Processing Unit} --- unidad central de procesamiento \\
CVE      & \textit{Common Vulnerabilities and Exposures} --- identificador público de una
           vulnerabilidad conocida \\
DDoS     & \textit{Distributed Denial of Service} --- denegación de servicio distribuida \\
DMZ      & \textit{Demilitarized Zone} --- zona desmilitarizada de una red \\
DORA     & \textit{Digital Operational Resilience Act} --- reglamento europeo de
           resiliencia operativa digital del sector financiero \\
EDR      & \textit{Endpoint Detection and Response} --- detección y respuesta en el
           \textit{endpoint} \\
FN       & Falso negativo \\
FP       & Falso positivo \\
GB       & Gigabyte \\
GPU      & \textit{Graphics Processing Unit} --- unidad de procesamiento gráfico \\
HIDS     & \textit{Host-based Intrusion Detection System} --- sistema de detección de
           intrusiones basado en \textit{host} \\
IoT      & \textit{Internet of Things} --- internet de las cosas \\
IP       & \textit{Internet Protocol} --- protocolo de internet \\
JSON     & \textit{JavaScript Object Notation} --- notación de objetos de JavaScript \\
LLM      & \textit{Large Language Model} --- modelo de lenguaje de gran escala \\
MDR      & \textit{Managed Detection and Response} --- detección y respuesta gestionada \\
NIS2     & \textit{Network and Information Security Directive 2} --- directiva europea de
           seguridad de redes y sistemas de información \\
NIST     & \textit{National Institute of Standards and Technology} --- instituto nacional
           de normas y tecnología de los Estados Unidos \\
ODF      & \textit{Optical Distribution Frame} --- repartidor óptico; punto de entrega del
           proveedor de servicio \\
OT       & \textit{Operational Technology} --- tecnología operativa \\
OWASP    & \textit{Open Web Application Security Project} --- proyecto abierto de
           seguridad en aplicaciones web \\
RAG      & \textit{Retrieval-Augmented Generation} --- generación aumentada por
           recuperación \\
RAM      & \textit{Random-Access Memory} --- memoria de acceso aleatorio \\
RF       & Requisito funcional \\
RNF      & Requisito no funcional \\
SCA      & \textit{Security Configuration Assessment} --- módulo de evaluación de
           configuración de seguridad de Wazuh \\
SIEM     & \textit{Security Information and Event Management} --- gestión de información
           y eventos de seguridad \\
SOAR     & \textit{Security Orchestration, Automation and Response} --- orquestación,
           automatización y respuesta de seguridad \\
SOC      & \textit{Security Operations Center} --- centro de operaciones de seguridad \\
SSH      & \textit{Secure Shell} --- protocolo de administración remota cifrada \\
TDP      & \textit{Thermal Design Power} --- potencia de diseño térmico \\
TR-069/TR-369 & \textit{Technical Report 069/369} --- protocolos del \textit{Broadband
           Forum} para la gestión remota de equipos de cliente \\
TUI      & \textit{Text-based User Interface} --- interfaz de usuario basada en texto \\
UCI      & \textit{Unified Configuration Interface} --- interfaz de configuración
           unificada de OpenWrt \\
USB      & Universidad Simón Bolívar \\
VN       & Verdadero negativo \\
VP       & Verdadero positivo \\
WSL      & \textit{Windows Subsystem for Linux} --- subsistema de Windows para Linux \\
XDR      & \textit{Extended Detection and Response} --- detección y respuesta extendida \\
YAML     & \textit{YAML Ain't Markup Language} --- formato de serialización de datos
           legible por personas \\
\end{longtable}
